\documentclass[11pt]{article}
\usepackage[a4paper,margin=2.6cm]{geometry}
\usepackage[T1]{fontenc}
\usepackage{lmodern}
\usepackage{amsmath,amssymb,amsthm}
\usepackage{booktabs}
\usepackage{enumitem}
\usepackage[hidelinks]{hyperref}

\newtheorem{theorem}{Theorem}
\newtheorem{lemma}[theorem]{Lemma}
\newtheorem{proposition}[theorem]{Proposition}
\newtheorem{corollary}[theorem]{Corollary}
\theoremstyle{definition}
\newtheorem{definition}[theorem]{Definition}
\newtheorem{assumption}{Assumption}
\renewcommand{\theassumption}{A\arabic{assumption}}
\theoremstyle{remark}
\newtheorem{remark}[theorem]{Remark}

\newcommand{\E}{\mathbb{E}}
\newcommand{\R}{\mathbb{R}}
\newcommand{\Prob}{\mathbb{P}}
\newcommand{\wga}{\mathrm{WGA}}
\newcommand{\code}[1]{\texttt{\small #1}}
\newcommand{\brk}{\allowbreak}
\emergencystretch=4em

\title{Concept Learning Curves and Early Concept Swap\\[2pt]
\large Single source of truth for the thesis method: idea, steps, theory, kill test and alternatives}
\author{Master's thesis working document, repository \code{Spur\_SpLiCE}}
\date{Version 1, 2026-11-14}

\begin{document}
\maketitle

\section*{How to read this document}

This file is the reference for every person and every coding model that works on the method. Section~\ref{sec:method}
specifies the method, Section~\ref{sec:theory} holds the theory with full proofs, Section~\ref{sec:kill} defines the
kill test on MetaShift, Section~\ref{sec:plan} lists the experiments, Section~\ref{sec:impl} gives implementation notes
for this repository and Section~\ref{sec:alt} lists three alternative topics.

Three conventions hold throughout.
\begin{itemize}[nosep]
\item \textbf{Status tags.} \textsc{[fixed]} marks a decision that implementers follow. \textsc{[default]} marks a
starting value that experiments may change. \textsc{[open]} marks a decision the student and the supervisor still take.
\item \textbf{Scope of the proofs.} Every theorem is proven inside the stylized model of Section~\ref{sec:model}. The
theorems show that the method is correct when Assumptions~\ref{as:sep} to~\ref{as:lin} hold. Whether real images
satisfy the assumptions is an empirical question, and Section~\ref{sec:kill} tests each assumption. A proof that the
method succeeds on real data does not exist for this method or for any published competitor; Section~\ref{sec:scope}
states the exact boundary.
\item \textbf{Language.} The confounding factor is the \emph{spurious attribute} or \emph{spurious concept}. Lists read
``A, B and C''.
\end{itemize}

\section{Problem and motivation}

A classifier trained on data where a spurious attribute $a$ correlates with the label $y$ relies on $a$ and fails on
the groups that break the correlation. Worst-group accuracy (WGA), the lowest accuracy over the groups $(y,a)$, measures
this failure. Networks pick up the feature that is easiest to fit first (simplicity bias), and on Waterbirds and
MetaShift that feature is the surroundings of the object.

The earlier project line, Spatial Concept Distillation, used patch-wise concept maps from CLIP and SpLiCE as regression
targets for the student. Its weakness is conceptual: the student copies CLIP, so the method is knowledge distillation
and the explanation tool explains the teacher. The present method keeps the patch-wise concept maps and changes their
role.

\begin{quote}
\textbf{Thesis statement.} Patch-wise concept maps turn the simplicity bias of the model under training into a
readable record: one learning curve per named concept. The concepts the model relies on first are flagged, and during
training their image regions are exchanged between images, which removes the correlation between those concepts and
the label. CLIP and SpLiCE supply names and locations. They supply no training target.
\end{quote}

The three components answer three questions.
\begin{itemize}[nosep]
\item \emph{Why SpLiCE?} The audit reports named concepts (``forest'', ``sofa''), which a human reads in seconds.
\item \emph{Why patches?} The intervention needs the location of a concept inside each image.
\item \emph{Why does it help on small datasets?} Methods that reweight examples (JTT, LfF, SPARE) depend on the few
images of the minority groups. A patch-level swap builds a counterfactual from every image
(Proposition~\ref{prop:ess}).
\end{itemize}

\section{Notation}

\begin{center}
\begin{tabular}{lp{0.68\linewidth}}
\toprule
Symbol & Meaning \\
\midrule
$x_i$, $y_i$, $a_i$ & image, class label, spurious attribute (hidden, evaluation only) of training example $i$ \\
$n$ & number of training images (MetaShift: about 1{,}700) \\
$K$ & number of concept units (single factors or groups of factors) \\
$M_i \in [0,1]^{K\times H\times W}$ & concept map of image $i$: per cell a distribution over the $K$ units, $H=W=14$ \\
$m_{ik}\in[0,1]^{H\times W}$ & region mask of unit $k$ in image $i$ (Definition~\ref{def:region}) \\
$\mathcal{I}_k$ & set of images in which unit $k$ is present \\
$f_t$ & model after $t$ epochs; $f_t(x)_y$ is the logit of class $y$ \\
$O_k$ & occlusion operator of unit $k$ (Definition~\ref{def:occ}) \\
$R_k(t)$, $S_k(t)$ & reliance of the model on unit $k$ at epoch $t$ and its share \\
$T_0$ & length of the early window in epochs \\
$\mathcal{F}$ & set of flagged concept units \\
$q$ & probability that a training image receives a swap \\
$p$, $\rho$ & $p=\Prob(a=y)$ and $\rho = 2p-1$, the correlation of $a$ and $y$ \\
\bottomrule
\end{tabular}
\end{center}

\section{The method, step by step}\label{sec:method}

The method has an offline stage (Step 0), an audit stage (Steps 1 to 3) and a training stage (Step 4).

\subsection*{Step 0: patch-wise concept maps \textsc{[fixed]}, existing code}

For every training image the frozen CLIP ViT-B/32 reads the image at 448 pixels through 224-pixel windows. The last
block uses query-query attention without the residual and feed-forward branch (ClearCLIP), which yields text-aligned
patch embeddings on a $14\times14$ grid. The SpLiCE concept dictionary, grouped into factors by the existing grouping
stage, supplies $K$ text directions. $M_i[k,h,w]$ is the softmax over the units of the cosine between the patch
embedding at $(h,w)$ and the direction of unit $k$. The maps are computed once and cached. No label enters.

\textsc{[default]} The concept unit is a \emph{group of factors that share pixels} (the 12 co-activation groups of
MetaShift obtained by spectral clustering of the factor maps), which keeps the audit cheap and its curves stable.
Single factors (130 on MetaShift) are the fine-grained variant.

\begin{definition}[Region and presence]\label{def:region}
The hard region of unit $k$ in image $i$ is the set of cells where $k$ has the largest map value. Unit $k$ is present in
image $i$ when its hard region holds at least three cells. The mask $m_{ik}$ is the hard region upsampled to the image
resolution with a soft border half a map cell wide.
\end{definition}

\subsection*{Step 1: occlusion and reliance \textsc{[fixed]}}

\begin{definition}[Occlusion operator]\label{def:occ}
$O_k(x_i)$ replaces the pixels of $x_i$ under $m_{ik}$ and leaves the other pixels unchanged. Two fillings exist.
\emph{Paint-over} fills the region by normalized convolution from its surroundings. \emph{Donor fill} takes the pixels
at the same positions from a random other training image. \textsc{[default]} The audit uses paint-over, which is
deterministic and already implemented.
\end{definition}

\begin{definition}[Reliance]\label{def:rel}
For a supervised model the reliance on unit $k$ at epoch $t$ is
\begin{equation}\label{eq:rel}
R_k(t) \;=\; \frac{1}{|\mathcal{I}_k|}\sum_{i\in\mathcal{I}_k}\Big(f_t(x_i)_{y_i}-f_t\big(O_k(x_i)\big)_{y_i}\Big)^2 ,
\end{equation}
and its share is $S_k(t)=R_k(t)/\sum_{j=1}^{K}R_j(t)$. For a self-supervised encoder $h_t$ the squared logit change
is replaced by $1-\cos\!\big(h_t(x_i),h_t(O_k(x_i))\big)$.
\end{definition}

Reliance is an explanation of the model under training: it measures how much the output of $f_t$ depends on the pixels
that carry concept $k$. CLIP appears only inside $m_{ik}$.

\textsc{[default]} Reliance is \emph{not divided by the region area}. Remark~\ref{rem:area} shows that area is part of
what makes a concept simple; the area-normalized score is an ablation. (This corrects the first oral description of the
idea, which proposed the normalized score as the main one.)

\subsection*{Step 2: concept learning curves \textsc{[fixed]}}

Train the model with plain ERM for $T_0$ epochs. After every epoch $t\le T_0$ compute $R_k(t)$ for all units on a fixed
probe subset of the training images, in evaluation mode and without augmentation. The $K$ sequences
$t\mapsto S_k(t)$ are the concept learning curves. \textsc{[default]} $T_0=20$ epochs and the probe subset is the full
MetaShift training set; larger datasets use 1{,}024 fixed images.

Cost per audit: one forward pass of the probe subset per unit, so $K+1$ passes. With $K=12$ and 1{,}700 images this is
below ten seconds on one V100.

\subsection*{Step 3: flagging \textsc{[default]}}

The onset score of unit $k$ is its mean early share,
\begin{equation}\label{eq:onset}
s_k \;=\; \frac{1}{T_0}\sum_{t=1}^{T_0} S_k(t).
\end{equation}
Sort the units by decreasing $s_k$. The flagged set $\mathcal{F}$ is the shortest prefix whose scores sum to at least
$\tau$, restricted so that the flagged regions cover at most 75 percent of the area of an average image.
\textsc{[default]} $\tau=0.5$.

The audit report lists every unit with its words, its onset score, its curve and its flag. The report is the xAI
deliverable of the thesis, and a human can overrule a flag before training continues (optional human-in-the-loop arm).

\subsection*{Step 4: early concept swap \textsc{[fixed]} structure, \textsc{[default]} values}

Restart training from the initial weights. For each training image $x_i$ in a batch, with probability $q$:
\begin{enumerate}[nosep]
\item compute the kept mask $\kappa_i = 1-\max_{k\in\mathcal{F}} m_{ik}$, the part of the image outside all flagged
regions;
\item draw a donor $x_j$ from the batch and paint over the donor's own kept-type content, so that the donor contributes
flagged-type content only;
\item form $\tilde x_i = \kappa_i\, x_i + (1-\kappa_i)\,\tilde x_j$ and train on $(\tilde x_i, y_i)$.
\end{enumerate}
The label stays $y_i$: the kept part decides the class. An image qualifies when its kept mask covers between 5 and 95
percent of the view. \textsc{[default]} $q=0.5$ for the whole run; the schedule ``$q$ high in the first third of
training, then decaying'' is the \emph{early-only} variant. The self-supervised variant uses $\tilde x_i$ as the second
view of $x_i$ in the SimCLR loss.

Step 4 reuses the existing compositing code with one change: the object mask of the earlier method, which came from
text prompts for objects and scenes, is replaced by the kept mask $\kappa_i$, which comes from the audit of the model.
The asymmetry ``keep this, exchange that'' now originates in the learning dynamics of the model.

\subsection*{Step 5: evaluation \textsc{[fixed]}}

Report test and validation WGA and average accuracy, per-group accuracy, the audit report and the detection quality of
Section~\ref{sec:kill}. Group labels are used for evaluation and for model selection at the level the baselines use
them.

\section{Theory}\label{sec:theory}

\subsection{Model and assumptions}\label{sec:model}

\begin{assumption}[Spatial separability]\label{as:sep}
Each image decomposes into concept regions, and the feature of a concept depends only on the pixels of its region. For
$K$ concepts the pooled feature is $z=(z_1,\dots,z_K)\in\R^K$, and occluding region $k$ replaces $z_k$ and leaves the
other coordinates unchanged.
\end{assumption}

\begin{assumption}[Two-concept Gaussian data]\label{as:data}
The label is $y\in\{-1,+1\}$ with $\Prob(y=1)=\tfrac12$. The spurious attribute is $a\in\{-1,+1\}$ with
$\Prob(a=y\mid y)=p\in(\tfrac12,1)$, so $\E[ya]=\rho=2p-1\in(0,1)$. There are a core concept and a spurious concept
with features
\[
z_c=\mu_c\,y+\sigma_c\,\xi_c,\qquad z_s=\mu_s\,a+\sigma_s\,\xi_s ,
\]
where $\mu_c,\mu_s,\sigma_c,\sigma_s>0$ and $\xi_c,\xi_s$ are independent standard normal variables, independent of
$(y,a)$. Write $m_c=\mu_c^2+\sigma_c^2$ and $m_s=\mu_s^2+\sigma_s^2$.
\end{assumption}

\begin{assumption}[Simplicity]\label{as:simple}
The spurious concept has the larger initial energy:
$\rho\,\mu_s\sqrt{m_s} \;>\; \mu_c\sqrt{m_c}$.
\end{assumption}

\begin{assumption}[Linear readout and squared loss]\label{as:lin}
The model is $f_w(z)=w^\top z$, trained by gradient flow on $L(w)=\tfrac12\,\E\big[(y-w^\top z)^2\big]$ from $w(0)=0$.
\end{assumption}

Under Assumption~\ref{as:data} a direct computation gives
\begin{equation}\label{eq:moments}
\Sigma=\E[zz^\top]=\begin{pmatrix} m_c & \rho\mu_c\mu_s\\ \rho\mu_c\mu_s & m_s\end{pmatrix},\qquad
b=\E[yz]=\begin{pmatrix}\mu_c\\ \rho\mu_s\end{pmatrix},
\end{equation}
using $\E[z_c]=\E[z_s]=0$, $\E[z_cz_s]=\mu_c\mu_s\E[ya]$ and $\E[yz_s]=\mu_s\E[ya]$. Since $\rho<1$,
$\det\Sigma=m_cm_s-\rho^2\mu_c^2\mu_s^2>0$.

\subsection{Reliance measures the weight of a concept}

\begin{lemma}[Occlusion reliance]\label{lem:occ}
Let $z\in\R^K$ have $\E[z_k]=0$ and $\E[z_k^2]=m_k$, and let $f_w(z)=w^\top z$. Under Assumption~\ref{as:sep}:
\begin{enumerate}[nosep,label=(\roman*)]
\item if occlusion sets $z_k$ to its mean $0$, then $R_k=\E[(f_w(z)-f_w(O_kz))^2]=w_k^2\,m_k$;
\item if occlusion replaces $z_k$ by an independent copy $z_k'$ with the same distribution, then $R_k=2\,w_k^2\,m_k$.
\end{enumerate}
In both cases the ranking of the concepts by $R_k$ equals their ranking by $|w_k|\sqrt{m_k}$.
\end{lemma}

\begin{proof}
By Assumption~\ref{as:sep}, $f_w(z)-f_w(O_kz)=w_k(z_k-\bar z_k)$, where $\bar z_k$ is the replacement. For (i),
$\bar z_k=0$ and $\E[w_k^2z_k^2]=w_k^2m_k$. For (ii), $\E[(z_k-z_k')^2]=\E[z_k^2]+\E[z_k'^2]-2\E[z_k]\E[z_k']=2m_k$.
The factor $2$ is common to all concepts.
\end{proof}

\subsection{The concept with the larger energy is learned first}

\begin{theorem}[Early dynamics]\label{thm:early}
Let $z\in\R^K$ with $\Sigma=\E[zz^\top]$, $b=\E[yz]$ and $\lambda=\lambda_{\max}(\Sigma)$. Under
Assumption~\ref{as:lin} the gradient flow satisfies
\[
w(t)=\int_0^t e^{-\Sigma s}\,b\,\mathrm{d}s
\qquad\text{and}\qquad
\big\|w(t)-t\,b\big\|_2\;\le\;\tfrac12\,\lambda\,\|b\|_2\,t^2\quad\text{for all }t\ge0 .
\]
\end{theorem}

\begin{proof}
$\nabla L(w)=\Sigma w-b$, so the flow is the linear equation $\dot w=b-\Sigma w$ with $w(0)=0$. Its unique solution is
$w(t)=\int_0^te^{-\Sigma s}b\,\mathrm{d}s$, as differentiation confirms:
$\dot w=e^{-\Sigma t}b$ and $b-\Sigma w(t)=b-\big(I-e^{-\Sigma t}\big)b=e^{-\Sigma t}b$. Hence
$w(t)-tb=\int_0^t\big(e^{-\Sigma s}-I\big)b\,\mathrm{d}s$. $\Sigma$ is symmetric positive semidefinite with
eigenvalues $\lambda_i\in[0,\lambda]$, so $I-e^{-\Sigma s}$ has eigenvalues $1-e^{-\lambda_is}\in[0,\lambda s]$
because $1-e^{-u}\le u$ for $u\ge0$. Therefore $\|(e^{-\Sigma s}-I)b\|_2\le\lambda s\|b\|_2$, and integrating over
$[0,t]$ gives the bound.
\end{proof}

\begin{corollary}[Order of the learning curves]\label{cor:order}
Define the initial energy of concept $k$ as $e_k=|b_k|\sqrt{m_k}$ with $m_k=\Sigma_{kk}$, and let reliance be given by
Lemma~\ref{lem:occ}. If $e_j>e_k$, then $R_j(t)>R_k(t)$ for every
\[
0<t<t^\star_{jk}=\frac{2\,(e_j-e_k)}{\lambda\,\|b\|_2\,\big(\sqrt{m_j}+\sqrt{m_k}\big)} .
\]
\end{corollary}

\begin{proof}
Let $\varepsilon=\tfrac12\lambda\|b\|_2t^2$. Theorem~\ref{thm:early} gives $|w_j(t)|\ge t|b_j|-\varepsilon$ and
$|w_k(t)|\le t|b_k|+\varepsilon$. Then
\[
|w_j(t)|\sqrt{m_j}-|w_k(t)|\sqrt{m_k}\;\ge\;t\,(e_j-e_k)-\varepsilon\big(\sqrt{m_j}+\sqrt{m_k}\big),
\]
which is positive exactly when $t<t^\star_{jk}$. By Lemma~\ref{lem:occ}, $R$ is an increasing function of
$|w|\sqrt{m}$ with the same constant for all concepts, so $R_j(t)>R_k(t)$.
\end{proof}

In the two-concept model, $e_s=\rho\mu_s\sqrt{m_s}$ and $e_c=\mu_c\sqrt{m_c}$ by \eqref{eq:moments}.
Assumption~\ref{as:simple} is $e_s>e_c$, so Corollary~\ref{cor:order} states: \emph{throughout the early window
$(0,t^\star_{sc})$ the spurious concept has the larger reliance, and the onset score \eqref{eq:onset} computed inside
this window ranks the spurious concept first.} This is the guarantee behind Steps 2 and 3.

\begin{remark}[Area is part of simplicity]\label{rem:area}
Let concept $k$ cover $n_k$ cells, each with signal $\nu_k$ and independent noise of variance $\tau_k^2$, and let the
feature be the sum over its cells. Then $\mu_k=n_k\nu_k$ and $\sigma_k^2=n_k\tau_k^2$, so
$e_k=|\mathrm{corr}_k|\,n_k\nu_k\sqrt{n_k^2\nu_k^2+n_k\tau_k^2}$ grows quadratically in the area $n_k$ once
$n_k\nu_k^2\ge\tau_k^2$. A large region is learned early because it is large. Dividing the reliance by the area would
remove a real part of the effect, which is why the unnormalized score is the default.
\end{remark}

\subsection{ERM keeps the spurious concept and loses worst-group accuracy}

\begin{proposition}[Limit of ERM]\label{prop:erm}
Under Assumptions~\ref{as:data} and~\ref{as:lin}, $w(t)\to w^\ast=\Sigma^{-1}b$ with
\[
w_c^\ast=\frac{\mu_c\big(\sigma_s^2+(1-\rho^2)\mu_s^2\big)}{\det\Sigma}>0,\qquad
w_s^\ast=\frac{\rho\,\mu_s\,\sigma_c^2}{\det\Sigma}>0,
\qquad
r^\ast(\rho)=\frac{w_s^\ast}{w_c^\ast}=\frac{\rho\,\mu_s\,\sigma_c^2}{\mu_c\big(\sigma_s^2+(1-\rho^2)\mu_s^2\big)} .
\]
\end{proposition}

\begin{proof}
$\Sigma$ is positive definite, so $e^{-\Sigma t}\to0$ and $w(t)=(I-e^{-\Sigma t})\Sigma^{-1}b\to\Sigma^{-1}b$. With
$\Sigma^{-1}=\frac{1}{\det\Sigma}\begin{pmatrix}m_s&-\rho\mu_c\mu_s\\-\rho\mu_c\mu_s&m_c\end{pmatrix}$ and
\eqref{eq:moments},
$w_c^\ast\det\Sigma=m_s\mu_c-\rho^2\mu_c\mu_s^2=\mu_c(\sigma_s^2+(1-\rho^2)\mu_s^2)$ and
$w_s^\ast\det\Sigma=m_c\rho\mu_s-\rho\mu_c^2\mu_s=\rho\mu_s\sigma_c^2$.
\end{proof}

The spurious weight is positive whenever the core feature is noisy ($\sigma_c>0$) and the correlation is positive: the
model stays reliant on the spurious concept at convergence.

\begin{theorem}[Worst-group accuracy decreases in the spurious ratio]\label{thm:wga}
Under Assumption~\ref{as:data}, classify by the sign of $w^\top z$ with $w_c>0$ and let $r=w_s/w_c$. The accuracy on
group $(y,a)$ is $\Phi\!\big((\mu_c+r\mu_s\,ya)/\sqrt{\sigma_c^2+r^2\sigma_s^2}\big)$, where $\Phi$ is the standard
normal distribution function, and
\[
\wga(r)=\Phi\big(g(|r|)\big),\qquad g(u)=\frac{\mu_c-u\,\mu_s}{\sqrt{\sigma_c^2+u^2\sigma_s^2}} .
\]
The function $g$ is strictly decreasing on $[0,\infty)$. Consequently $\wga$ is maximal at $r=0$ with value
$\Phi(\mu_c/\sigma_c)$, and it decreases strictly as $|r|$ grows.
\end{theorem}

\begin{proof}
Given $(y,a)$, $y\,w^\top z=w_c\mu_c+w_s\mu_s\,ya+w_c\sigma_cy\xi_c+w_s\sigma_sy\xi_s$ is normal with mean
$w_c(\mu_c+r\mu_s\,ya)$ and variance $w_c^2(\sigma_c^2+r^2\sigma_s^2)$. The prediction is correct when this variable
is positive, which has probability $\Phi(\text{mean}/\text{standard deviation})$; dividing by $w_c>0$ gives the group
accuracy. The groups with $ya=-\mathrm{sign}(r)$ have the smaller mean $\mu_c-|r|\mu_s$ at equal variance, so they
form the worst group and $\wga(r)=\Phi(g(|r|))$. For $u\ge0$,
\[
g'(u)=\frac{-\mu_s(\sigma_c^2+u^2\sigma_s^2)-(\mu_c-u\mu_s)\,u\,\sigma_s^2}{(\sigma_c^2+u^2\sigma_s^2)^{3/2}}
=\frac{-\mu_s\sigma_c^2-\mu_c\,u\,\sigma_s^2}{(\sigma_c^2+u^2\sigma_s^2)^{3/2}}<0 .
\]
$\Phi$ is strictly increasing, which completes the proof.
\end{proof}

Proposition~\ref{prop:erm} and Theorem~\ref{thm:wga} together give
$\wga_{\mathrm{ERM}}=\Phi(g(r^\ast(\rho)))<\Phi(\mu_c/\sigma_c)$: ERM is strictly below the best worst-group accuracy
a linear classifier reaches.

\subsection{Swapping the flagged concept restores worst-group accuracy}

\begin{definition}[Swap in the model]\label{def:swap}
A swap of the spurious concept replaces $z_s$ by the spurious feature $z_s'$ of an independent training example and
keeps $(y,z_c)$. Each training example is swapped independently with probability $q\in[0,1]$.
\end{definition}

\begin{theorem}[Early concept swap]\label{thm:swap}
Under Assumptions~\ref{as:sep}, \ref{as:data} and~\ref{as:lin}, train on the distribution obtained by swapping the
spurious concept with probability $q$. Then:
\begin{enumerate}[nosep,label=(\roman*)]
\item the swapped distribution is the distribution of Assumption~\ref{as:data} with the correlation
$\rho_q=(1-q)\rho$ and all other parameters unchanged;
\item the trained predictor has the spurious ratio $r^\ast(\rho_q)$, which is strictly decreasing in $q$ with
$r^\ast(\rho_1)=0$;
\item the worst-group accuracy on the \emph{original} test distribution, $\Phi(g(r^\ast(\rho_q)))$, is strictly
increasing in $q$ and reaches $\Phi(\mu_c/\sigma_c)$ at $q=1$, the maximum over all linear classifiers with $w_c>0$.
\end{enumerate}
\end{theorem}

\begin{proof}
(i) The donor attribute $a'$ is independent of $(y,\xi_c)$, and $\Prob(a'=1)=\tfrac12$ because the classes are
balanced and the attribute is symmetric. The attribute after the swap equals $a$ with probability $1-q$ and $a'$ with
probability $q$, so it equals $y$ with probability $(1-q)p+\tfrac q2$, independently of the noise. The correlation is
$2\big((1-q)p+\tfrac q2\big)-1=(1-q)(2p-1)=(1-q)\rho$. The donor noise is standard normal and independent, so the
feature after the swap has the form $\mu_s\tilde a+\sigma_s\tilde\xi_s$ required by Assumption~\ref{as:data}.

(ii) By (i) and Proposition~\ref{prop:erm} the limit has ratio
$r^\ast(\rho_q)=\rho_q\mu_s\sigma_c^2/\big(\mu_c(\sigma_s^2+(1-\rho_q^2)\mu_s^2)\big)$. The numerator increases
strictly in $\rho_q$ and the denominator decreases in $\rho_q$, so $r^\ast$ increases strictly in $\rho_q$ and
decreases strictly in $q$. At $q=1$, $\rho_q=0$ and $r^\ast=0$.

(iii) The test distribution is unchanged by the training intervention, so Theorem~\ref{thm:wga} applies with
$r=r^\ast(\rho_q)\ge0$. Since $g$ is strictly decreasing and $r^\ast(\rho_q)$ is strictly decreasing in $q$, the
composition is strictly increasing in $q$, and the value at $r=0$ is the maximum by Theorem~\ref{thm:wga}.
\end{proof}

\begin{remark}[Imperfect masks]\label{rem:imperfect}
If the flagged regions cover the spurious concept in a fraction $\pi$ of the swapped images, the effective swap
probability is $q\pi$ and Theorem~\ref{thm:swap} holds with $\rho_{q\pi}$. The gain degrades smoothly with the mask
quality and stays positive for every $\pi>0$.
\end{remark}

\subsection{Why the swap needs the audit: an untyped swap can lower WGA}

\begin{proposition}[Swapping both concepts]\label{prop:sym}
Under Assumptions~\ref{as:sep}, \ref{as:data} and~\ref{as:lin}, let each training example have its core concept
swapped with probability $q_c$, its spurious concept swapped with probability $q_s$ (disjoint events,
$q_c+q_s\le1$ and $q_c<1$), and keep its label. With $\alpha=1-q_c$, $\beta=1-q_s$ and $\gamma=1-q_c-q_s$, the trained predictor
has the ratio
\[
r(q_c,q_s)=\frac{\rho\,\mu_s\,\big(\beta\,m_c-\alpha\gamma\,\mu_c^2\big)}
{\mu_c\,\big(\alpha\,m_s-\beta\gamma\,\rho^2\mu_s^2\big)} .
\]
For the symmetric swap $q_c=q_s=q/2$,
\[
r_{\mathrm{sym}}(q)=\frac{\rho\,\mu_s\,\big(\sigma_c^2+q\,\mu_c^2\big)}
{\mu_c\,\big(\sigma_s^2+\mu_s^2(1-(1-q)\rho^2)\big)},
\]
which is strictly increasing in $q$ if and only if
$\mu_c^2\big(\sigma_s^2+(1-\rho^2)\mu_s^2\big)>\rho^2\mu_s^2\sigma_c^2$, that is, if and only if
$w_c^\ast\mu_c>w_s^\ast\rho\mu_s$. In that regime the symmetric swap lowers the worst-group accuracy below ERM.
\end{proposition}

\begin{proof}
On a core swap, $z_c$ is independent of $(y,z_s)$; on a spurious swap, $z_s$ is independent of $(y,z_c)$. The marginal
distributions of $z_c$ and $z_s$ are unchanged, so the diagonal of $\Sigma$ is unchanged. Averaging over the three
events gives $\E[yz_c]=\alpha\mu_c$, $\E[yz_s]=\beta\rho\mu_s$ and $\E[z_cz_s]=\gamma\rho\mu_c\mu_s$. Solving
$\Sigma w=b$ as in Proposition~\ref{prop:erm},
$w_s\propto m_c\beta\rho\mu_s-\gamma\rho\mu_c\mu_s\cdot\alpha\mu_c$ and
$w_c\propto m_s\alpha\mu_c-\gamma\rho\mu_c\mu_s\cdot\beta\rho\mu_s$ with the same positive factor, which gives
$r(q_c,q_s)$. For $q_c=q_s=q/2$ we have $\alpha=\beta$ and $\gamma=1-q$, so
$m_c-\gamma\mu_c^2=\sigma_c^2+q\mu_c^2$ and $m_s-\gamma\rho^2\mu_s^2=\sigma_s^2+\mu_s^2(1-(1-q)\rho^2)$. Writing
$N_0=\sigma_c^2$ and $D_0=\sigma_s^2+(1-\rho^2)\mu_s^2$, $r_{\mathrm{sym}}(q)$ is a positive constant times
$(N_0+q\mu_c^2)/(D_0+q\rho^2\mu_s^2)$, whose derivative has the sign of $\mu_c^2D_0-\rho^2\mu_s^2N_0$. By
Proposition~\ref{prop:erm} this is the sign of $w_c^\ast\mu_c-w_s^\ast\rho\mu_s$. The last claim follows from
Theorem~\ref{thm:wga}.
\end{proof}

Proposition~\ref{prop:sym} explains why the swap must be typed. Exchanging every concept alike weakens the core signal
as much as the spurious one and raises the relative weight of the spurious concept whenever the core concept carries
the larger share of the ERM margin. The earlier experiments with symmetric concept mixing, which tied their rectangle
control on MetaShift and lost to spatial maps on Waterbirds, agree with this result. The audit of Steps 1 to 3 supplies
the type without a text prior.

\subsection{Why a swap suits small datasets}

\begin{proposition}[Effective sample size]\label{prop:ess}
Consider $n$ examples of which a fraction $p$ belongs to the majority groups. Reweighting that gives the majority and
the minority groups equal total weight uses the weights $1/(2p)$ and $1/(2(1-p))$ and has the Kish effective sample
size
\[
n_{\mathrm{eff}}=\frac{(\sum_iw_i)^2}{\sum_iw_i^2}=4\,p\,(1-p)\,n .
\]
A swap with $q=1$ gives every example the weight $1$ and has $n_{\mathrm{eff}}=n$.
\end{proposition}

\begin{proof}
$\sum_iw_i=n\big(p\cdot\tfrac{1}{2p}+(1-p)\cdot\tfrac{1}{2(1-p)}\big)=n$ and
$\sum_iw_i^2=n\big(\tfrac{1}{4p}+\tfrac{1}{4(1-p)}\big)=\tfrac{n}{4p(1-p)}$.
\end{proof}

For a weighted mean of independent observations with common variance $v$, the variance is $v/n_{\mathrm{eff}}$. Both
routes decorrelate $a$ and $y$ in expectation; the swap does so at the full sample size. MetaShift with $p=0.88$ keeps
$42$ percent of its sample under reweighting (about 720 of 1{,}700 images) and Waterbirds with $p=0.95$ keeps $19$
percent. Reweighting by inferred groups (JTT, LfF, SPARE) inherits this bound at best.

\subsection{What the theory establishes and where it stops}\label{sec:scope}

Inside the model the chain is complete: the spurious concept leads the learning curves in an explicit time window
(Corollary~\ref{cor:order}), reliance reads this order from occlusions (Lemma~\ref{lem:occ}), ERM converges to a
predictor with a positive spurious weight and a suboptimal WGA (Proposition~\ref{prop:erm}, Theorem~\ref{thm:wga}),
and the swap of the flagged concept raises WGA monotonically to the optimum (Theorem~\ref{thm:swap}).

Four gaps separate the model from a ResNet on photographs. Each is an assumption with an empirical test.

\begin{center}
\begin{tabular}{p{0.20\linewidth}p{0.40\linewidth}p{0.30\linewidth}}
\toprule
Assumption & Risk on real data & Test \\
\midrule
\ref{as:sep} separability & CLIP maps at $14\times14$ are coarse; a receptive field of the student covers several
regions; hair colour and gender (CelebA) share pixels & map quality check and mask-quality ablation
(Remark~\ref{rem:imperfect}) \\
\ref{as:data} data model & more than two concepts, non-Gaussian features, donor fills leave seams & random-region
control with equal area \\
\ref{as:simple} simplicity & the core concept may have the larger energy; then the audit flags the core concept and
Proposition~\ref{prop:sym} with $q_s=0$ predicts harm & kill test, criterion K1 \\
\ref{as:lin} linear model & deep networks learn features; a shortcut learned early can persist & restart-versus-continue
ablation; curves on the real network \\
\bottomrule
\end{tabular}
\end{center}

Assumption~\ref{as:simple} cannot be derived. From the training distribution alone, the core and the spurious concept
play symmetric roles: both are features correlated with the label. Every method without group labels adds an assumption
that breaks this symmetry, and ``the spurious feature is the one learned first'' is the assumption shared with LfF, JTT
and SPARE. The thesis states it, measures how often it holds and reports the audit so a reader can verify the flags.

Theorem~\ref{thm:early} extends to wide networks in the kernel regime, where the same linear dynamics hold in function
space with the neural tangent kernel in place of $\Sigma$, and deep linear networks learn the directions of largest
singular value first (Saxe et al., 2014). The cross-entropy loss has the same gradient at $w=0$ up to a constant
factor, so the early ordering of Corollary~\ref{cor:order} carries over. These extensions are cited, and their proofs
lie outside this document.

For the self-supervised variant the theoretical support is the identifiability result of von K\"ugelgen et al.\ (2021):
contrastive learning recovers the block of latent factors that the augmentation leaves unchanged. A swap of the flagged
regions moves the flagged concepts into the block that changes between the two views.

\section{Kill test on MetaShift}\label{sec:kill}

MetaShift (cat against dog, indoor against outdoor, about 1{,}700 training images, 88 percent aligned) is the
development dataset: one run takes minutes. The test has two stages and stops at the first failure.

\subsection*{Stage K1: does the audit find the spurious concept? (about one week)}

\begin{enumerate}[nosep]
\item Train ERM for $T_0=20$ epochs, seeds 1 to 4, and store $R_k(t)$ for every unit and epoch.
\item Type every concept unit with the hidden labels through the existing diagnostic
(\code{cospro.\brk diagnostics.\brk factor\_\brk validity}): a unit is an \emph{attribute} unit or a \emph{class} unit when its
conditional information score reaches 0.05 and doubles the other score.
\item Compute the AUROC of the onset score $s_k$ for separating attribute units from class units, and the Spearman
correlation of the unit ranking between seeds.
\end{enumerate}

\textbf{Pass:} AUROC at least 0.7 on three of four seeds, the flagged set contains scene units (lawn, desks, beds) and
excludes the cat and dog units, and the mean rank correlation between seeds reaches 0.6.
\textbf{Fail:} AUROC near 0.5, or the animal units lead the curves. A failure ends the timing hypothesis; the fallback
is the human-in-the-loop arm (a person ticks the spurious units in the audit report) or an alternative of
Section~\ref{sec:alt}.

Run K1 for two settings in parallel, since the outcome decides the thesis setting \textsc{[open]}: ResNet-18 from
scratch and ImageNet-pretrained ResNet-50.

\subsection*{Stage K2: does the swap raise WGA? (about two weeks)}

Arms, four seeds each, all with the same schedule and model selection:
\begin{center}
\begin{tabular}{ll}
\toprule
Arm & Purpose \\
\midrule
ERM & reference \\
swap, audit flags ($\mathcal{F}$) & the method \\
swap, random units of equal area & tests whether the flagged identity matters \\
swap, last-ranked units & tests the direction of the timing signal \\
swap, shuffled maps & tests whether the concept maps matter \\
swap, oracle scene units & upper bound with the correct type \\
\bottomrule
\end{tabular}
\end{center}

\textbf{Pass:} the method beats ERM and the random-units control on validation WGA on at least three of four seeds,
with a mean gain of at least 5 points over ERM, and the order ``oracle $\ge$ method $>$ random $>$ last-ranked'' holds
on average. MetaShift validation groups hold 72 images, so single-seed differences below 7 points carry no weight.
\textbf{Fail:} the method ties its random control. Then the gain comes from mixing as such, and the thesis moves to an
alternative.

\section{Experiment plan after the kill test}\label{sec:plan}

\begin{itemize}[nosep]
\item \textbf{Datasets.} MetaShift (development), Waterbirds (locked test with 642 images per minority group),
Spur-CIFAR10 and one further dataset with large evaluation groups. CelebA violates Assumption~\ref{as:sep} and serves
as the stated limit.
\item \textbf{Baselines.} ERM, JTT, LfF or SPARE, MaskTune, Decompose-and-Compose, DFR as the reference with group
labels, and object-prior compositing from the earlier project.
\item \textbf{Ablations.} Area-normalized reliance; single factors against groups as units; paint-over against donor
fill in the audit; $T_0\in\{5,10,20,40\}$; $\tau\in\{0.3,0.5,0.7\}$; $q\in\{0.25,0.5,0.75,1\}$; restart against
continue after the audit; early-only schedule; human-corrected flags.
\item \textbf{Theory checks on real data.} WGA as a function of $q$ (Theorem~\ref{thm:swap} predicts a monotone rise);
the spurious ratio measured by reliance at the end of training as a function of $q$; a synthetic two-concept dataset
that satisfies the assumptions exactly, where the formulas of Section~\ref{sec:theory} are verified numerically.
\item \textbf{Timeline (14 weeks).} Weeks 1 to 3 kill test; weeks 4 to 6 method, controls and ablations on MetaShift
and Waterbirds; weeks 7 to 9 baselines, seeds and further datasets; weeks 10 to 14 writing.
\end{itemize}

\section{Implementation notes for this repository}\label{sec:impl}

These notes are for coding models. Names of existing functions come from
\code{docs/\brk notes/\brk CONCEPT\_\brk FACTORS\_\brk 2026\_\brk 09\_\brk 26.\brk md}; confirm each name in the code before use.

\begin{itemize}[nosep]
\item \textbf{Concept maps.} \code{cospro.\brk cli.\brk build\_\brk concept\_\brk maps} with launcher
\code{scripts/\brk build\_\brk concept\_\brk maps.\brk sbatch} writes the $14\times14$ maps; MetaShift maps exist. Co-activation groups
come from \code{coactivation\_\brk groups}.
\item \textbf{Warping.} Training views are random crops and flips. The existing \code{warp\_\brk maps} maps the cached
concept maps onto a view from its crop box and flip. The audit uses unaugmented images.
\item \textbf{Swap.} Extend the compositing path (\code{TrainingMethod.\brk compose\_\brk views}, \code{composite\_\brk maps}): the
object mass becomes the kept mask $\kappa_i$ built from $\mathcal{F}$. Keep the sharpening of the mask and the
paint-over of the donor.
\item \textbf{New code.} (a) A supervised ERM training method with a class head, mixed precision
(\code{torch.\brk float16} autocast with a \code{GradScaler}) and W\&B logging of WGA at every evaluation; (b) a reliance
callback that computes \eqref{eq:rel} per unit and epoch and writes \code{concept\_\brk curves.\brk json} under \code{outputs/\brk };
(c) a flagging CLI that reads the curves and writes \code{flagged\_\brk units.\brk json} and the audit report; (d) the option
that feeds the flagged units into the swap.
\item \textbf{Cluster rules.} All training runs through \code{sbatch}. Every new entry point ships with a launcher in
\code{scripts/\brk } that starts from \code{scripts/\brk \_\brk template.\brk sbatch}, sources
\code{scripts/\brk load\_\brk splice\_\brk cluster\_\brk env.\brk sh} and calls \code{announce\_\brk results}. One V100, at most 5 CPUs and 8G of
memory per CPU. At most 8 jobs run at once (\code{MAX\_\brk PARALLEL}).
\item \textbf{Storage.} Result JSON under \code{outputs/\brk }; checkpoints and tensors under scratch through
\code{cospro.\brk tracking.\brk artifacts}, written once at the end of a stage.
\item \textbf{Tests.} Unit tests for the reliance callback on a synthetic two-concept fixture where the true weights
are known, and for the flagging rule. \code{tests/\brk test\_\brk slurm\_\brk launchers.\brk py} covers the launchers.
\item \textbf{Numerical check of the theory.} A short script that samples the model of Assumption~\ref{as:data},
solves $\Sigma w=b$ under the swap and compares the ratio and the WGA with Proposition~\ref{prop:erm},
Theorem~\ref{thm:swap} and Proposition~\ref{prop:sym}.
\end{itemize}

\section{Open decisions}

\begin{enumerate}[nosep]
\item \textsc{[open]} \textbf{Supervised or self-supervised as the main setting.} The theory, the reliance on a class
logit and the published baselines belong to the supervised setting, which is the recommendation. The self-supervised
variant reuses the existing SimCLR pipeline and can remain a secondary table.
\item \textsc{[open]} \textbf{Backbone.} From scratch (stronger simplicity bias, lower absolute numbers) or
ImageNet-pretrained ResNet-50 (the protocol of the literature). Stage K1 runs both.
\item \textsc{[open]} \textbf{Model selection.} Validation WGA with group labels, as JTT and Decompose-and-Compose use
it, or validation average accuracy for a fully label-free claim.
\end{enumerate}

\section{Position against prior work}

A literature search on 2026-11-14 (about 16 queries) found no work that combines named patch-level concepts, the
timing of their reliance as the signal for spuriousness and a region-level intervention. The search was not exhaustive;
repeat it before submission. The closest works are listed with their difference.

\begin{center}
\begin{tabular}{p{0.27\linewidth}p{0.30\linewidth}p{0.33\linewidth}}
\toprule
Work & Shared element & Difference \\
\midrule
DISC (Wu et al., 2023) & concept sensitivity at every epoch & image-level concepts, a concept image bank and
environments \\
Decompose-and-Compose (Noohdani et al., 2024) & attribution-based parts recombined across images & unnamed parts;
group-labelled validation selects the causal part \\
MaskTune (Asgari et al., 2022) & masks what the model relies on & after training, one epoch, masks all salient input \\
JTT, LfF, SPARE & early-training signal & example-level reweighting \\
GALS (Petryk et al., 2022) & CLIP attention guides a classifier & fixed class prompt as the prior \\
RaVL (Varma et al., 2024) & region-level discovery & fine-tuned vision-language models, region clusters \\
SpLiCE (Bhalla et al., 2024) & sparse named concepts & image-level, explains CLIP itself \\
\bottomrule
\end{tabular}
\end{center}

\section{Alternative thesis topics}\label{sec:alt}

Each alternative reuses the cached patch-wise concept maps. Their novelty has had no dedicated literature search yet.

\subsection*{Alternative A: counterfactual last-layer retraining}

\textbf{Idea.} Deep feature reweighting (DFR) retrains the last layer on a small group-balanced set and needs group
labels for it. Build that set synthetically: composite every training image with exchanged concept regions so that
concept and class are independent, extract features with the frozen ERM backbone and retrain the last layer on them.

\textbf{Why it may work.} Theorem~\ref{thm:swap} applies to the last layer exactly, since the last layer is a linear
model on fixed features; Assumption~\ref{as:lin} holds by construction. Training takes seconds, so MetaShift allows
hundreds of seeds and configurations.

\textbf{Risk.} DFR works when the backbone already encodes the core feature; composites with seams may shift the
features. Related work: DFR, Decompose-and-Compose.
\textbf{Effort.} Lowest of all options, about four weeks to a full result table. It also serves as a fast first test of
the swap before full training.

\subsection*{Alternative B: concept-localized attribution penalty}

\textbf{Idea.} Keep the audit of Steps 1 to 3 and replace the swap by a penalty on the model's attribution: the
input-gradient or Grad-CAM mass inside the flagged regions is penalized during training, in the manner of ``right for
the right reasons''.

\textbf{Why it may work.} No composite images, hence no seams and no donor statistics. In the linear model the penalty
is a ridge term on $w_s$ alone, which shrinks the ratio $r$ monotonically, so Theorem~\ref{thm:wga} gives the gain.

\textbf{Risk.} GALS uses a CLIP attention map with a class prompt for the same penalty; the difference is the
model-derived named flag set. Input-gradient penalties cost a second backward pass.
\textbf{Effort.} Moderate; shares the audit code with the main method.

\subsection*{Alternative C: patch-concept bottleneck classifier}

\textbf{Idea.} Classify from the patch-wise SpLiCE codes directly: each patch is a sparse vector over named concepts,
an attention pooling aggregates the patches and a linear layer predicts the class. The model is interpretable by
construction, and removing a concept is an edit of one weight column. The audit decides which concepts to remove.

\textbf{Why it may work.} The backbone is frozen CLIP, so 1{,}700 images suffice, and concept removal in a linear
bottleneck is exact.

\textbf{Risk.} Closest prior work is strong: spatially grounded concept bottlenecks report 87 percent WGA on
Waterbirds (arXiv 2510.04180), and image-level concept removal on CLIP is established. The contribution would be the
patch-wise SpLiCE codes and the label-free choice of the removed concepts.
\textbf{Effort.} Low to moderate; no training of a backbone.

\subsection*{Minimal fallback}

If no intervention yields a gain, the concept learning curves remain a diagnostic contribution: a study of which named
concepts ERM, JTT, DFR and the self-supervised methods rely on over training, on four datasets.

\begin{thebibliography}{99}\small
\bibitem{asgari} Asgari, Khani, Khasahmadi, Aminbeidokhti et al. MaskTune: Mitigating Spurious Correlations by Forcing
to Explore. NeurIPS 2022.
\bibitem{bhalla} Bhalla, Oesterling, Srinivas, Calmon and Lakkaraju. Interpreting CLIP with Sparse Linear Concept
Embeddings (SpLiCE). NeurIPS 2024.
\bibitem{kalimeris} Kalimeris et al. SGD on Neural Networks Learns Functions of Increasing Complexity. NeurIPS 2019.
\bibitem{kirichenko} Kirichenko, Izmailov and Wilson. Last Layer Re-Training is Sufficient for Robustness to Spurious
Correlations. ICLR 2023.
\bibitem{lan} Lan et al. ClearCLIP: Decomposing CLIP Representations for Dense Vision-Language Inference. ECCV 2024.
\bibitem{liu} Liu et al. Just Train Twice: Improving Group Robustness without Training Group Information. ICML 2021.
\bibitem{nam} Nam, Cha, Ahn, Lee and Shin. Learning from Failure: Training Debiased Classifier from Biased Classifier.
NeurIPS 2020.
\bibitem{noohdani} Noohdani et al. Decompose-and-Compose: A Compositional Approach to Mitigating Spurious Correlation.
CVPR 2024.
\bibitem{petryk} Petryk et al. On Guiding Visual Attention with Language Specification. CVPR 2022.
\bibitem{ross} Ross, Hughes and Doshi-Velez. Right for the Right Reasons. IJCAI 2017.
\bibitem{sagawa} Sagawa, Raghunathan, Koh and Liang. An Investigation of Why Overparameterization Exacerbates Spurious
Correlations. ICML 2020.
\bibitem{saxe} Saxe, McClelland and Ganguli. Exact Solutions to the Nonlinear Dynamics of Learning in Deep Linear
Neural Networks. ICLR 2014.
\bibitem{shah} Shah, Tamuly, Raghunathan, Jain and Netrapalli. The Pitfalls of Simplicity Bias in Neural Networks.
NeurIPS 2020.
\bibitem{varma} Varma et al. RaVL: Discovering and Mitigating Spurious Correlations in Fine-Tuned Vision-Language
Models. NeurIPS 2024.
\bibitem{vonk} von K\"ugelgen et al. Self-Supervised Learning with Data Augmentations Provably Isolates Content from
Style. NeurIPS 2021.
\bibitem{wu} Wu, Yuksekgonul, Zhang and Zou. Discover and Cure: Concept-aware Mitigation of Spurious Correlation.
ICML 2023.
\bibitem{yang} Yang, Gan, Dziugaite and Mirzasoleiman. Identifying Spurious Biases Early in Training through the Lens
of Simplicity Bias. AISTATS 2024.
\end{thebibliography}

\end{document}
